\section{结论}
\label{sec:conclusion}

本次审计的三项核心发现重新界定了物理约束风电预测的实证基础。

在文献报道层面，46 项被审计文献（2023--2026）均未披露任何失效工况、负面权衡或运行边界（0.0\%）。鉴于 95.7\% 的研究缺乏公开代码、50.0\% 将神经模块切除混同为物理消融，文献中一致正向的叙事表明现有评估与报道协议未对边界工况展开有效压力测试。

在软件实现层面，在被审计的 46 项研究中仅有的两个公开代码库中，PINN-UQ 经审查证实代码对齐，而作为端到端物理约束 Transformer 唯一公开训练代码库的 PCT，其执行的损失函数颠倒了 Jensen--Shannon 散度的参数方向，这一执行偏差在不审查源码的情况下无法察觉。代数推导独立证实其核乘积在数学上保持双变量耦合，表明实现偏差与模型有效性位于不同的计算层次。

在运行决策层面，物理先验的实际决策价值取决于样本规模、气象代理时效对齐以及下游市场结算规则。在 $N_{\text{train}}=1{,}583$ 且 $\lambda=1.0$ 条件下，真实功率曲线先验将归一化调度成本降低 $1.28\%$；完全相同条件下形式与权重匹配的伪打乱先验将其抬升 $61.2\%$，揭示出模型在强约束下对先验误设的高度敏感性。在弱约束尺度（$\lambda=0.01$）下，真实先验与伪打乱先验表现接近，实证更符合非特异性数值平滑；当训练损失中物理代理时效错配超过 $k > 12$ 步时，三座风场的调度成本均出现单调恶化。

PCT 代码审计、零申报现金流异常与打乱先验激增并非三个孤立问题，共享同一根源：该领域所采用的评价指标与消融协议无法区分物理知识与数值正则化、实现保真度与论文声明、预测技术与市场结算结构。覆盖约束溯源、软件执行、代理陈旧度与决策解耦评估的完整审计，虽较常规基准测试耗费更多工作量，但回答了对于电网运行决策更具现实意义的方法学问题。

在我们测试的配置下，经验功率曲线先验在受限训练窗口下提供了可测量的归纳价值，但其运维收益需要严谨的实证边界。界定这些条件需要参数对齐的消融对照、打乱阴性控制实验、静态计算图验证，以及将远期合同收益与实时不平衡结算现金流分离的市场评估协议。
